\subsection{WLAN-Anbindung eingebetteter Systeme}
\gls{wlan}-fähige Mikrocontroller wie der ESP32 können in realen Anwendungsumgebungen eingesetzt werden, beispielsweise in Innenräumen wie Wohnungen, Büros oder öffentlichen Bereichen. 
Durch die Nutzung vorhandener Funkkommunikation lassen sich Systeme ohne zusätzliche spezialisierte Infrastruktur realisieren. 
Dies ermöglicht eine einfache Integration in bestehende technische Umgebungen~\cite{hernandez2022wifi}.
Die Einbindung des Systems in ein lokales Netzwerk erfolgt über die integrierte \gls{wlan}-Schnittstelle des Mikrocontrollers~\cite{WaveshareESP32C6}.

Für die Kommunikation in Netzwerken ist eine eindeutige Adressierung der beteiligten Geräte erforderlich. 
Diese erfolgt über \gls{ip}-Adressen, die jedem Teilnehmer zugewiesen werden und den gezielten Datenaustausch ermöglichen. 
Auch eingebettete Systeme benötigen eine solche Adresse, um innerhalb eines \gls{wlan} kommunizieren zu können.
Die Vergabe dieser Adressen kann automatisiert durch das \gls{dhcp} erfolgen. 
Dabei werden Geräten beim Verbinden mit einem Netzwerk selbstständig \gls{ip}-Adressen sowie weitere notwendige Konfigurationsparameter zugewiesen. 
Dies reduziert den manuellen Konfigurationsaufwand und erleichtert die Verwaltung von Netzwerken erheblich~\cite{dhcp2023}.

Nachdem das Gerät eine \gls{ip}-Adresse  zugewiesen bekommen hat, kann es über diese im Netzwerk angesprochen werden.
Durch den Einsatz von \gls{dhcp} können Geräte ohne manuelle Netzwerkkonfiguration in ein bestehendes Netzwerk integriert werden. 
Dies ist insbesondere bei eingebetteten Systemen von Vorteil, da diese häufig in großer Anzahl eingesetzt werden und eine individuelle Konfiguration mit hohem Aufwand verbunden wäre.
Insbesondere in \gls{wlan}-Umgebungen ermöglicht \gls{dhcp} eine einfache und schnelle Einbindung neuer Geräte. 
Dies zeigt sich sowohl in privaten Netzwerken als auch in öffentlichen Zugangspunkten, bei denen sich Geräte ohne zusätzliche Konfigurationsschritte verbinden und unmittelbar kommunizieren können~\cite{dhcp2023}.